\documentclass[11pt]{article}
\usepackage[landscape,margin=0.8in]{geometry}
\usepackage{amsmath,amssymb}
\usepackage{graphicx}
\usepackage{booktabs}
\usepackage{caption}
\usepackage{float}
\usepackage{xcolor}
% The auto-generated \input tables still carry some Japanese header labels,
% so we compile with xelatex + xeCJK. English body text is unaffected.
\usepackage[hidelinks]{hyperref}
\usepackage{etoolbox}
% Start every top-level section on a new page
\pretocmd{\section}{\clearpage}{}{}

\newcommand{\code}[1]{\texttt{#1}}
\newcommand{\kt}{\kappa_t}
\newcommand{\Nbar}{\bar N}
\newcommand{\Nhat}{\hat N}
\graphicspath{{../results/figures/}}

\title{\textbf{Market Concentration and the Flattening of the Phillips Curve}\\[2pt]
\large A Theory-Motivated, Semi-Structural State-Space Assessment of the HSA Mechanism}
\author{Satoshi Ichikawa}
\date{Draft --- August 2026}

\begin{document}
% ================================================================
\section{Research Questions}
% ================================================================

The empirical analysis is organized around a semi-structural state-space
representation of the HSA mechanism. The objective is not to estimate the
complete structural model or to recover the underlying HSA demand system.
Instead, the analysis asks whether the central empirical implications of the
HSA Phillips curve are supported by U.S. aggregate data.

Throughout this section, the competition measure is oriented so that a larger
value of $N_t$ represents a more competitive environment, or equivalently a
larger effective number of firms. Hence, a decline in $N_t$ corresponds to
greater market concentration. If a concentration measure such as an HHI is
used directly instead, the directional signs stated below must be reversed.

Let the observed competition measure admit the state-space representation

\begin{equation}
N_t^{\mathrm{obs}}
=
\bar N_t
+
\hat N_t
+
\nu_t^{N},
\label{eq:competition_measurement}
\end{equation}

where $\bar N_t$ denotes a slow-moving competitive environment,
$\hat N_t$ denotes a stationary latent deviation around that environment,
and $\nu_t^{N}$ denotes measurement error.

For interpretation and reporting, define

\begin{equation}
\tilde N_t
=
\bar N_t-N_0,
\label{eq:competition_center}
\end{equation}

where $N_0$ is a fixed reporting center.

The empirical Phillips curve is represented as

\begin{equation}
\pi_t
=
\alpha E_t\pi_{t+1}
+
\kappa_t x_t
+
\theta_t \hat N_t
+
\varepsilon_t,
\label{eq:empirical_nkpc}
\end{equation}

where $x_t$ denotes the Phillips-curve forcing variable. The two
competition-dependent coefficients are parameterized locally as

\begin{equation}
\kappa_t
=
\kappa_0+\delta\tilde N_t,
\label{eq:kappa_state}
\end{equation}

and

\begin{equation}
\theta_t
=
\theta_0+\gamma\tilde N_t.
\label{eq:theta_state}
\end{equation}

The coefficients $\kappa_t$ and $\theta_t$ are empirical varying
coefficients. Their structural interpretation is conditional on the
maintained mapping between the observed competition measure, the latent
competition states, and the competitive environment in the HSA model.

The directional implications below are not imposed mechanically in the
baseline estimation. The parameters are estimated without sign restrictions,
and the posterior probability assigned to the theory-consistent region is
reported as an empirical measure of mechanism consistency.

% ------------------------------------------------
\subsection{RQ1: Competition-Dependent Phillips-Curve Slope}
% ------------------------------------------------

\subsubsection*{Does the conditional response of inflation to the forcing variable vary systematically with the competitive environment?}

The first research question concerns the slow-moving component of the
Phillips-curve slope. Equation~\eqref{eq:kappa_state} allows the conditional
slope to vary with the competitive environment:

\begin{equation}
\kappa_t
=
\kappa_0+\delta\tilde N_t.
\end{equation}

The parameter $\delta$ measures how the conditional Phillips-curve slope
changes with the slow-moving competition state.

Under the maintained HSA interpretation, greater market concentration
reduces the structural slope of the Phillips curve. Since larger values of
$N_t$ are defined here to represent greater competition, the directional
hypothesis is

\begin{equation}
H_0^{(1)}:\delta\leq0
\qquad\text{against}\qquad
H_1^{(1)}:\delta>0.
\label{eq:hyp_delta}
\end{equation}

Thus, a positive $\delta$ implies that the Phillips curve becomes steeper
when competition increases and flatter when market concentration rises.

In the Bayesian analysis, the corresponding posterior quantity of interest is

\begin{equation}
\Pr\!\left(
\delta>0
\mid
\mathcal D
\right),
\end{equation}

where $\mathcal D$ denotes the observed data.

% ------------------------------------------------
\subsection{RQ2: Short-Run Direct Competition Channel}
% ------------------------------------------------

\subsubsection*{Does the stationary competition state contain inflation-relevant information beyond the standard Phillips-curve forcing variable?}

The second research question concerns the direct loading of inflation on the
stationary competition state. Holding the interaction in
Equation~\eqref{eq:theta_state} fixed at the reporting center gives

\begin{equation}
\pi_t
=
\alpha E_t\pi_{t+1}
+
\kappa_t x_t
+
\theta_0\hat N_t
+
\varepsilon_t.
\end{equation}

The coefficient $\theta_0$ measures the direct short-run association between
the stationary competition state and inflation after conditioning on expected
inflation and the Phillips-curve forcing variable.

In the HSA mechanism, endogenous movements in the number of firms generate a
cost-push component distinct from the conventional marginal-cost term. With
$N_t$ oriented so that larger values correspond to greater competition, the
theory-consistent direct loading is negative. The directional hypothesis is
therefore

\begin{equation}
H_0^{(2)}:\theta_0\geq0
\qquad\text{against}\qquad
H_1^{(2)}:\theta_0<0.
\label{eq:hyp_theta}
\end{equation}

The corresponding posterior quantity is

\begin{equation}
\Pr\!\left(
\theta_0<0
\mid
\mathcal D
\right).
\end{equation}

In the baseline empirical interpretation, $\theta_0$ is treated
conservatively as a direct latent-state loading. Its interpretation as the
endogenous-entry or participation channel of the HSA model is conditional on
the maintained theoretical mapping.

% ------------------------------------------------
\subsection{RQ3: Competition-Dependent Dynamic Effect}
% ------------------------------------------------

\subsubsection*{Does the strength of the short-run competition channel itself vary with the slow-moving competitive environment?}

The third research question asks whether the direct competition loading is
itself state dependent:

\begin{equation}
\theta_t
=
\theta_0+\gamma\tilde N_t.
\end{equation}

The coefficient $\gamma$ measures how the short-run effect associated with
$\hat N_t$ changes as the economy moves across different competitive
environments.

Under the maintained HSA interpretation with the relevant pass-through
channel, greater concentration strengthens the magnitude of the negative
endogenous competition effect. Since lower values of $N_t$ represent greater
concentration, this implies that $\theta_t$ should become more negative as
$\bar N_t$ falls. In the linear approximation above, the resulting
directional hypothesis is

\begin{equation}
H_0^{(3)}:\gamma\leq0
\qquad\text{against}\qquad
H_1^{(3)}:\gamma>0.
\label{eq:hyp_gamma}
\end{equation}

A positive $\gamma$ therefore implies

\begin{equation}
\bar N_t\downarrow
\quad\Longrightarrow\quad
\theta_t\downarrow,
\end{equation}

so that the direct short-run competition channel becomes more negative as
market concentration increases.

The corresponding posterior quantity is

\begin{equation}
\Pr\!\left(
\gamma>0
\mid
\mathcal D
\right).
\end{equation}

This prediction is stronger than the existence of the direct effect itself.
The sign of $\theta_0$ is motivated by the endogenous cost-push channel,
whereas the directional prediction for $\gamma$ additionally requires that
the magnitude of that channel varies systematically with the competitive
environment.

% ------------------------------------------------
\subsection{RQ4: Joint HSA Mechanism Consistency}
% ------------------------------------------------

\subsubsection*{Are the slope and dynamic competition channels jointly organized in the direction implied by the HSA mechanism?}

The fourth research question evaluates the mechanism jointly rather than
testing each coefficient in isolation. The empirical Phillips curve contains
two conceptually distinct competition channels:

\begin{equation}
\underbrace{\kappa_t x_t}_{\text{competition-dependent slope channel}}
\qquad\text{and}\qquad
\underbrace{\theta_t\hat N_t}_{\text{direct dynamic competition channel}}.
\end{equation}

The first channel captures systematic variation in the conditional
Phillips-curve slope across competitive environments. The second captures a
direct inflation contribution from short-run competition dynamics that is not
absorbed by the conventional forcing variable.

For the orientation of $N_t$ maintained above, the full theory-consistent
sign configuration is

\begin{equation}
\delta>0,
\qquad
\theta_0<0,
\qquad
\gamma>0.
\label{eq:hsa_sign_region}
\end{equation}

Rather than reducing this multidimensional restriction to a single
coefficient test, the Bayesian analysis evaluates the joint posterior
probability

\begin{equation}
\Pr\!\left(
\delta>0,\,
\theta_0<0,\,
\gamma>0
\mid
\mathcal D
\right).
\label{eq:joint_hsa_probability}
\end{equation}

This quantity measures how much posterior probability is assigned to the
region in which all three directional implications of the empirical HSA
mechanism hold simultaneously.

The analysis also reports the three marginal directional probabilities

\begin{equation}
\Pr(\delta>0\mid\mathcal D),
\qquad
\Pr(\theta_0<0\mid\mathcal D),
\qquad
\Pr(\gamma>0\mid\mathcal D),
\end{equation}

so that a failure of joint mechanism consistency can be traced to a
particular empirical channel.

% ------------------------------------------------
\subsection{Nested Model Interpretation}
% ------------------------------------------------

The research questions generate a nested sequence of empirical
specifications:

\begin{align}
\text{Model 1:}\qquad
& \delta=\theta_0=\gamma=0,
\\
\text{Model 2:}\qquad
& \delta\ \text{free},
\qquad
\theta_0=\gamma=0,
\\
\text{Model 3:}\qquad
& \delta,\theta_0\ \text{free},
\qquad
\gamma=0,
\\
\text{Model 4:}\qquad
& \delta,\theta_0,\gamma\ \text{free}.
\end{align}

Model 1 is the competition-invariant benchmark. It allows a conventional
Phillips-curve relationship but rules out both competition channels.

Model 2 introduces the competition-dependent slope channel and therefore
addresses RQ1.

Model 3 additionally introduces the direct stationary competition channel
and therefore addresses RQ2 while retaining the slope channel from Model 2.

Model 4 allows the direct competition coefficient itself to vary with the
slow-moving competitive environment and therefore addresses RQ3 and the full
joint mechanism in RQ4.

The sequence can therefore be summarized as

\begin{equation}
\text{Competition-Invariant}
\;\longrightarrow\;
\text{Slope Channel}
\;\longrightarrow\;
\text{Slope + Direct Channel}
\;\longrightarrow\;
\text{Full State-Dependent HSA}.
\end{equation}

Model comparison is based on posterior parameter distributions, marginal
likelihoods, and predictive performance. The purpose is not only to determine
whether a time-varying Phillips curve fits the data better, but also to
determine whether its time variation is systematically linked to the
competitive environment in the direction implied by the HSA mechanism.

% ------------------------------------------------
\subsection{Mechanism Decomposition}
% ------------------------------------------------

Conditional on the estimated state paths, Equation~\eqref{eq:empirical_nkpc}
permits the competition-related contribution to inflation to be separated
into two components.

The slope channel is

\begin{equation}
C_t^{\mathrm{slope}}
=
\left(
\kappa_t-\kappa_0
\right)x_t
=
\delta\tilde N_t x_t,
\end{equation}

while the direct dynamic channel is

\begin{equation}
C_t^{\mathrm{dynamic}}
=
\theta_t\hat N_t
=
\left(
\theta_0+\gamma\tilde N_t
\right)\hat N_t.
\end{equation}

Their sum,

\begin{equation}
C_t^{\mathrm{HSA}}
=
C_t^{\mathrm{slope}}
+
C_t^{\mathrm{dynamic}},
\end{equation}

provides an empirical decomposition of the competition-related component of
fitted inflation.

This decomposition distinguishes a change in the conditional transmission
of the forcing variable from a separate short-run competition contribution.
It therefore provides the empirical basis for assessing whether the observed
flattening is primarily associated with a changing Phillips-curve slope, a
distinct dynamic competition channel, or both.

% ================================================================
\section{Case 1: Quarterly Joint Estimation with Capital IQ Competition Data}
% ================================================================

Capital IQ provides a quarterly competition measure, so inflation and
competition are jointly observed at the quarterly frequency. The competition
series is decomposed into a slow-moving component and a stationary cyclical
component.

\begin{equation}
\mathbf{s}_t
=
\begin{bmatrix}
\tilde N_t \\
\hat N_t
\end{bmatrix},
\qquad
\tilde N_t=\bar N_t-N_0,
\qquad
n_t^{\mathrm{CIQ}}=N_t^{\mathrm{CIQ}}-N_0.
\end{equation}

The transition equation is

\begin{equation}
\mathbf{s}_t
=
\mathbf F\mathbf{s}_{t-1}
+
\boldsymbol{\eta}_t,
\qquad
\mathbf F
=
\begin{bmatrix}
1 & 0 \\
0 & \rho_N
\end{bmatrix},
\qquad
|\rho_N|<1.
\end{equation}

\begin{equation}
\boldsymbol{\eta}_t
\sim
\mathcal N(\mathbf 0,\mathbf Q),
\qquad
\mathbf Q
=
\begin{bmatrix}
\sigma_{\bar N}^2 & 0 \\
0 & \sigma_{\hat N}^2
\end{bmatrix}.
\end{equation}

The joint observation vector and measurement disturbances are

\begin{equation}
\mathbf y_t
=
\begin{bmatrix}
\pi_t \\
n_t^{\mathrm{CIQ}}
\end{bmatrix},
\qquad
\boldsymbol{\varepsilon}_t
\sim
\mathcal N(\mathbf 0,\mathbf R),
\qquad
\mathbf R
=
\begin{bmatrix}
\sigma_\pi^2 & 0 \\
0 & \sigma_{\nu,N}^2
\end{bmatrix}.
\end{equation}

For compact notation, define

\begin{equation}
\mathbf d_t
=
\begin{bmatrix}
\alpha E_t\pi_{t+1}+\kappa_0x_t \\
0
\end{bmatrix}.
\end{equation}


% ------------------------------------------------
\subsection{Model 0: CES}
% ------------------------------------------------

The CES benchmark imposes a constant Phillips-curve slope and excludes any
direct competition channel.

\begin{equation}
\mathbf y_t
=
\mathbf d_t
+
\begin{bmatrix}
0 & 0 \\
1 & 1
\end{bmatrix}
\mathbf s_t
+
\boldsymbol{\varepsilon}_t,
\qquad
\delta=\theta_0=\gamma=0.
\end{equation}


% ------------------------------------------------
\subsection{Model 1: HSA Competition Effect on Slope}
% ------------------------------------------------

Model 1 allows the slow-moving competitive environment to change the
Phillips-curve slope.

\begin{equation}
\kappa_t=\kappa_0+\delta\tilde N_t,
\qquad
\mathbf y_t
=
\mathbf d_t
+
\begin{bmatrix}
\delta x_t & 0 \\
1 & 1
\end{bmatrix}
\mathbf s_t
+
\boldsymbol{\varepsilon}_t.
\end{equation}

The specification remains conditionally linear and Gaussian because $x_t$
is observed.


% ------------------------------------------------
\subsection{Model 2: HSA Direct Competition Channel}
% ------------------------------------------------

Model 2 keeps the Phillips-curve slope constant but allows the stationary
competition state to affect inflation directly.

\begin{equation}
\mathbf y_t
=
\mathbf d_t
+
\begin{bmatrix}
0 & \theta_0 \\
1 & 1
\end{bmatrix}
\mathbf s_t
+
\boldsymbol{\varepsilon}_t,
\qquad
\delta=\gamma=0.
\end{equation}


% ------------------------------------------------
\subsection{Model 3: HSA Competition Dynamic Effect}
% ------------------------------------------------

Model 3 allows the loading on the stationary competition state to depend on
the slow-moving competitive environment.

\begin{equation}
\theta_t=\theta_0+\gamma\tilde N_t,
\qquad
\delta=0.
\end{equation}

\begin{equation}
\mathbf y_t
=
\mathbf d_t
+
\mathbf h_3(\mathbf s_t)
+
\boldsymbol{\varepsilon}_t,
\qquad
\mathbf h_3(\mathbf s_t)
=
\begin{bmatrix}
\theta_0\hat N_t+\gamma\tilde N_t\hat N_t \\
\tilde N_t+\hat N_t
\end{bmatrix}.
\end{equation}

Because the inflation equation contains the product
$\tilde N_t\hat N_t$, the measurement equation is bilinear in the latent
states.


% ------------------------------------------------
\subsection{Model 4: HSA Joint Competition Effect on Slope and Direct Competition Channel}
% ------------------------------------------------

Model 4 combines the competition-dependent slope and the state-dependent
direct competition channel.

\begin{equation}
\kappa_t=\kappa_0+\delta\tilde N_t,
\qquad
\theta_t=\theta_0+\gamma\tilde N_t.
\end{equation}

\begin{equation}
\mathbf y_t
=
\mathbf d_t
+
\mathbf h_4(\mathbf s_t)
+
\boldsymbol{\varepsilon}_t,
\qquad
\mathbf h_4(\mathbf s_t)
=
\begin{bmatrix}
\delta x_t\tilde N_t
+\theta_0\hat N_t
+\gamma\tilde N_t\hat N_t \\
\tilde N_t+\hat N_t
\end{bmatrix}.
\end{equation}

The competition-related contribution to inflation can be decomposed as

\begin{equation}
C_t^{\mathrm{slope}}=\delta\tilde N_t x_t,
\qquad
C_t^{\mathrm{direct}}=\theta_0\hat N_t,
\qquad
C_t^{\mathrm{interaction}}=\gamma\tilde N_t\hat N_t.
\end{equation}

\begin{equation}
C_t^{\mathrm{dynamic}}
=
C_t^{\mathrm{direct}}+C_t^{\mathrm{interaction}},
\qquad
C_t^{\mathrm{HSA}}
=
C_t^{\mathrm{slope}}+C_t^{\mathrm{dynamic}}.
\end{equation}



% ================================================================
\section{Case 2: Quarterly Joint Estimation with Interpolated Gustavo Competition Data}
% ================================================================

The annual Gustavo competition series is first interpolated to the quarterly
frequency. The resulting quarterly series is treated as the observed
competition measure in the joint state-space system.

\begin{equation}
\mathbf{s}_t
=
\begin{bmatrix}
\tilde N_t \\
\hat N_t
\end{bmatrix},
\qquad
\tilde N_t=\bar N_t-N_0,
\qquad
n_t^{\mathrm{G,PCHIP}}
=
N_t^{\mathrm{G,PCHIP}}-N_0.
\end{equation}

\begin{equation}
\mathbf{s}_t
=
\mathbf F\mathbf{s}_{t-1}
+
\boldsymbol{\eta}_t,
\qquad
\mathbf F
=
\begin{bmatrix}
1 & 0 \\
0 & \rho_N
\end{bmatrix},
\qquad
|\rho_N|<1.
\end{equation}

\begin{equation}
\boldsymbol{\eta}_t
\sim
\mathcal N(\mathbf 0,\mathbf Q),
\qquad
\mathbf Q
=
\begin{bmatrix}
\sigma_{\bar N}^2 & 0 \\
0 & \sigma_{\hat N}^2
\end{bmatrix}.
\end{equation}

\begin{equation}
\mathbf y_t
=
\begin{bmatrix}
\pi_t \\
n_t^{\mathrm{G,PCHIP}}
\end{bmatrix},
\qquad
\mathbf d_t
=
\begin{bmatrix}
\alpha E_t\pi_{t+1}+\kappa_0x_t \\
0
\end{bmatrix}.
\end{equation}

\begin{equation}
\boldsymbol{\varepsilon}_t
\sim
\mathcal N(\mathbf 0,\mathbf R),
\qquad
\mathbf R
=
\begin{bmatrix}
\sigma_\pi^2 & 0 \\
0 & \sigma_{\nu,N}^2
\end{bmatrix}.
\end{equation}


% ------------------------------------------------
\subsection{Model 0: CES}
% ------------------------------------------------

The CES benchmark excludes both HSA competition channels.

\begin{equation}
\mathbf y_t
=
\mathbf d_t
+
\begin{bmatrix}
0 & 0 \\
1 & 1
\end{bmatrix}
\mathbf s_t
+
\boldsymbol{\varepsilon}_t,
\qquad
\delta=\theta_0=\gamma=0.
\end{equation}


% ------------------------------------------------
\subsection{Model 1: HSA Competition Effect on Slope}
% ------------------------------------------------

The slow competition state shifts the conditional Phillips-curve slope.

\begin{equation}
\kappa_t=\kappa_0+\delta\tilde N_t,
\qquad
\mathbf y_t
=
\mathbf d_t
+
\begin{bmatrix}
\delta x_t & 0 \\
1 & 1
\end{bmatrix}
\mathbf s_t
+
\boldsymbol{\varepsilon}_t.
\end{equation}


% ------------------------------------------------
\subsection{Model 2: HSA Direct Competition Channel}
% ------------------------------------------------

The stationary competition state enters inflation directly.

\begin{equation}
\mathbf y_t
=
\mathbf d_t
+
\begin{bmatrix}
0 & \theta_0 \\
1 & 1
\end{bmatrix}
\mathbf s_t
+
\boldsymbol{\varepsilon}_t,
\qquad
\delta=\gamma=0.
\end{equation}


% ------------------------------------------------
\subsection{Model 3: HSA Competition Dynamic Effect}
% ------------------------------------------------

The direct competition loading varies with the slow competitive environment.

\begin{equation}
\theta_t=\theta_0+\gamma\tilde N_t,
\qquad
\delta=0.
\end{equation}

\begin{equation}
\mathbf y_t
=
\mathbf d_t
+
\begin{bmatrix}
\theta_0\hat N_t+\gamma\tilde N_t\hat N_t \\
\tilde N_t+\hat N_t
\end{bmatrix}
+
\boldsymbol{\varepsilon}_t.
\end{equation}


% ------------------------------------------------
\subsection{Model 4: HSA Joint Competition Effect on Slope and Direct Competition Channel}
% ------------------------------------------------

The full HSA specification allows both competition channels to operate
simultaneously.

\begin{equation}
\kappa_t=\kappa_0+\delta\tilde N_t,
\qquad
\theta_t=\theta_0+\gamma\tilde N_t.
\end{equation}

\begin{equation}
\mathbf y_t
=
\mathbf d_t
+
\begin{bmatrix}
\delta x_t\tilde N_t
+\theta_0\hat N_t
+\gamma\tilde N_t\hat N_t \\
\tilde N_t+\hat N_t
\end{bmatrix}
+
\boldsymbol{\varepsilon}_t.
\end{equation}

\begin{equation}
C_t^{\mathrm{slope}}=\delta\tilde N_t x_t,
\qquad
C_t^{\mathrm{direct}}=\theta_0\hat N_t,
\qquad
C_t^{\mathrm{interaction}}=\gamma\tilde N_t\hat N_t.
\end{equation}



% ================================================================
\section{Case 3: Mixed-Frequency Joint Estimation with Gustavo Competition Data}
% ================================================================

The annual Gustavo competition measure is not interpolated. Instead, the
observed annual value is assigned to Q4, while Q1--Q3 are left missing.
The underlying competition process is nevertheless defined at the quarterly
frequency and estimated jointly within the state-space model.

\begin{equation}
\mathbf s_t
=
\begin{bmatrix}
\tilde N_t \\
\hat N_t
\end{bmatrix},
\qquad
\tilde N_t=\bar N_t-N_0.
\end{equation}

The quarterly latent competition states evolve according to

\begin{equation}
\mathbf s_t
=
\mathbf F\mathbf s_{t-1}
+
\boldsymbol{\eta}_t,
\qquad
\mathbf F
=
\begin{bmatrix}
1 & 0 \\
0 & \rho_N
\end{bmatrix},
\qquad
|\rho_N|<1.
\end{equation}

\begin{equation}
\boldsymbol{\eta}_t
\sim
\mathcal N(\mathbf 0,\mathbf Q),
\qquad
\mathbf Q
=
\begin{bmatrix}
\sigma_{\bar N}^{2} & 0 \\
0 & \sigma_{\hat N}^{2}
\end{bmatrix}.
\end{equation}

The observed annual competition series is

\begin{equation}
n_t^{\mathrm{G,obs}}
=
\begin{cases}
N_t^{\mathrm G}-N_0,
& t\in\mathcal Q_4, \\[4pt]
\text{missing},
& t\in\mathcal Q_1\cup\mathcal Q_2\cup\mathcal Q_3.
\end{cases}
\end{equation}

Its measurement equation is

\begin{equation}
n_t^{\mathrm{G,obs}}
=
\tilde N_t+\hat N_t+\nu_t^N,
\qquad
\nu_t^N\sim\mathcal N(0,\sigma_{\nu,N}^{2}),
\qquad
t\in\mathcal Q_4.
\end{equation}

For Q1--Q3, the competition observation is missing and therefore contributes
no measurement density. The quarterly latent states continue to evolve
through the transition equation and are inferred jointly from the complete
state-space system.

The inflation equation is observed every quarter. Define

\begin{equation}
d_t
=
\alpha E_t\pi_{t+1}+\kappa_0x_t,
\qquad
\varepsilon_t^\pi\sim\mathcal N(0,\sigma_\pi^2).
\end{equation}


% ------------------------------------------------
\subsection{Model 0: CES}
% ------------------------------------------------

The CES benchmark excludes competition from the inflation equation.

\begin{equation}
\pi_t
=
d_t+\varepsilon_t^\pi,
\qquad
\delta=\theta_0=\gamma=0.
\end{equation}

Together with the competition measurement equation,

\begin{equation}
\begin{bmatrix}
\pi_t \\
n_t^{\mathrm{G,obs}}
\end{bmatrix}
=
\begin{bmatrix}
d_t \\
\tilde N_t+\hat N_t
\end{bmatrix}
+
\begin{bmatrix}
\varepsilon_t^\pi \\
\nu_t^N
\end{bmatrix},
\qquad
t\in\mathcal Q_4.
\end{equation}

For Q1--Q3 only the first measurement equation is observed.


% ------------------------------------------------
\subsection{Model 1: HSA Competition Effect on Slope}
% ------------------------------------------------

The slow-moving competition state changes the Phillips-curve slope.

\begin{equation}
\kappa_t=\kappa_0+\delta\tilde N_t,
\qquad
\pi_t
=
d_t+\delta x_t\tilde N_t+\varepsilon_t^\pi.
\end{equation}

Thus, in Q4,

\begin{equation}
\begin{bmatrix}
\pi_t \\
n_t^{\mathrm{G,obs}}
\end{bmatrix}
=
\begin{bmatrix}
d_t+\delta x_t\tilde N_t \\
\tilde N_t+\hat N_t
\end{bmatrix}
+
\begin{bmatrix}
\varepsilon_t^\pi \\
\nu_t^N
\end{bmatrix}.
\end{equation}

In Q1--Q3, $n_t^{\mathrm{G,obs}}$ remains missing while the latent quarterly
competition states are propagated by the transition equation.


% ------------------------------------------------
\subsection{Model 2: HSA Direct Competition Channel}
% ------------------------------------------------

The stationary competition state enters inflation directly.

\begin{equation}
\pi_t
=
d_t+\theta_0\hat N_t+\varepsilon_t^\pi,
\qquad
\delta=\gamma=0.
\end{equation}

In Q4,

\begin{equation}
\begin{bmatrix}
\pi_t \\
n_t^{\mathrm{G,obs}}
\end{bmatrix}
=
\begin{bmatrix}
d_t+\theta_0\hat N_t \\
\tilde N_t+\hat N_t
\end{bmatrix}
+
\begin{bmatrix}
\varepsilon_t^\pi \\
\nu_t^N
\end{bmatrix}.
\end{equation}

In Q1--Q3, quarterly inflation helps update the latent competition state
through its direct loading even though the annual competition observation
is missing.


% ------------------------------------------------
\subsection{Model 3: HSA Competition Dynamic Effect}
% ------------------------------------------------

The direct loading itself varies with the slow-moving competition state.

\begin{equation}
\theta_t=\theta_0+\gamma\tilde N_t,
\qquad
\pi_t
=
d_t
+
\theta_0\hat N_t
+
\gamma\tilde N_t\hat N_t
+
\varepsilon_t^\pi.
\end{equation}

In Q4,

\begin{equation}
\begin{bmatrix}
\pi_t \\
n_t^{\mathrm{G,obs}}
\end{bmatrix}
=
\begin{bmatrix}
d_t+\theta_0\hat N_t+\gamma\tilde N_t\hat N_t \\
\tilde N_t+\hat N_t
\end{bmatrix}
+
\begin{bmatrix}
\varepsilon_t^\pi \\
\nu_t^N
\end{bmatrix}.
\end{equation}

Because $\tilde N_t\hat N_t$ enters the inflation equation, this specification
has a bilinear measurement equation.


% ------------------------------------------------
\subsection{Model 4: HSA Joint Competition Effect on Slope and Direct Competition Channel}
% ------------------------------------------------

The full specification combines the slope and dynamic competition channels.

\begin{equation}
\kappa_t=\kappa_0+\delta\tilde N_t,
\qquad
\theta_t=\theta_0+\gamma\tilde N_t.
\end{equation}

\begin{equation}
\pi_t
=
d_t
+
\delta x_t\tilde N_t
+
\theta_0\hat N_t
+
\gamma\tilde N_t\hat N_t
+
\varepsilon_t^\pi.
\end{equation}

In Q4, the complete measurement system is

\begin{equation}
\begin{bmatrix}
\pi_t \\
n_t^{\mathrm{G,obs}}
\end{bmatrix}
=
\begin{bmatrix}
d_t
+\delta x_t\tilde N_t
+\theta_0\hat N_t
+\gamma\tilde N_t\hat N_t
\\
\tilde N_t+\hat N_t
\end{bmatrix}
+
\begin{bmatrix}
\varepsilon_t^\pi \\
\nu_t^N
\end{bmatrix}.
\end{equation}

In Q1--Q3,

\begin{equation}
n_t^{\mathrm{G,obs}}=\text{missing},
\end{equation}

so only the quarterly inflation measurement is used. The latent competition
states are not interpolated beforehand; their quarterly paths are inferred
inside the state-space model from the transition dynamics, the Q4 competition
observations, and the quarterly inflation equation.

\begin{equation}
C_t^{\mathrm{slope}}=\delta\tilde N_t x_t,
\qquad
C_t^{\mathrm{direct}}=\theta_0\hat N_t,
\qquad
C_t^{\mathrm{interaction}}=\gamma\tilde N_t\hat N_t.
\end{equation}

% ================================================================
\section{Case 4: Mixed-Frequency Joint Estimation with Gustavo Competition Data and Number of Establishments}
% ================================================================

The annual Gustavo competition measure is observed only in Q4 and is not
interpolated. Quarterly movements between the annual observations are inferred
within the state-space model using both the latent competition dynamics and
quarterly changes in the number of establishments.

Let $E_t$ denote the quarterly number of establishments and define the
demeaned quarterly establishment-growth signal as

\begin{equation}
g_t^{E}
=
\Delta\log E_t
-
\overline{\Delta\log E}.
\end{equation}

The establishment series is used as an auxiliary signal for short-run
business-demography movements rather than as a direct measure of competition.

The latent competition state is

\begin{equation}
\mathbf s_t
=
\begin{bmatrix}
\tilde N_t \\
\hat N_t
\end{bmatrix},
\qquad
\tilde N_t=\bar N_t-N_0.
\end{equation}

The quarterly transition equation is

\begin{equation}
\mathbf s_t
=
\mathbf F\mathbf s_{t-1}
+
\mathbf B g_t^{E}
+
\boldsymbol{\eta}_t,
\qquad
\mathbf F
=
\begin{bmatrix}
1 & 0 \\
0 & \rho_N
\end{bmatrix},
\qquad
\mathbf B
=
\begin{bmatrix}
0 \\
\lambda_E
\end{bmatrix}.
\end{equation}

Equivalently,

\begin{equation}
\tilde N_t
=
\tilde N_{t-1}+\eta_t^{\bar N},
\qquad
\hat N_t
=
\rho_N\hat N_{t-1}
+
\lambda_E g_t^{E}
+
\eta_t^{\hat N},
\qquad
|\rho_N|<1.
\end{equation}

The state innovations satisfy

\begin{equation}
\boldsymbol{\eta}_t
\sim
\mathcal N(\mathbf 0,\mathbf Q),
\qquad
\mathbf Q
=
\begin{bmatrix}
\sigma_{\bar N}^{2} & 0 \\
0 & \sigma_{\hat N}^{2}
\end{bmatrix}.
\end{equation}

Thus, $\rho_N$ captures persistence in short-run competition dynamics, while
$\lambda_E$ determines how quarterly establishment fluctuations inform the
latent competition state.

The annual Gustavo observation is retained only in Q4:

\begin{equation}
n_t^{\mathrm{G,obs}}
=
\begin{cases}
N_t^{\mathrm G}-N_0,
& t\in\mathcal Q_4, \\[4pt]
\text{missing},
& t\in\mathcal Q_1\cup\mathcal Q_2\cup\mathcal Q_3.
\end{cases}
\end{equation}

At every observed Q4, the latent states are required to reproduce the Gustavo
competition measure exactly:

\begin{equation}
n_t^{\mathrm{G,obs}}
=
\begin{bmatrix}
1 & 1
\end{bmatrix}
\mathbf s_t
=
\tilde N_t+\hat N_t,
\qquad
t\in\mathcal Q_4.
\end{equation}

Hence, the annual observations act as exact anchors for the quarterly latent
competition path. In Q1--Q3, no artificial competition observation is created;
the missing quarterly values are inferred from the transition dynamics,
establishment fluctuations, and the inflation equation.

For compact notation, define

\begin{equation}
d_t
=
\alpha E_t\pi_{t+1}
+
\kappa_0x_t,
\qquad
\varepsilon_t^\pi
\sim
\mathcal N(0,\sigma_\pi^2).
\end{equation}


% ------------------------------------------------
\subsection{Model 0: CES}
% ------------------------------------------------

The CES benchmark uses establishment growth only to recover the latent
competition path; competition does not enter the inflation equation.

\begin{equation}
\pi_t
=
d_t+\varepsilon_t^\pi,
\qquad
\delta=\theta_0=\gamma=0.
\end{equation}

In Q4, the joint measurement system is

\begin{equation}
\begin{bmatrix}
\pi_t \\
n_t^{\mathrm{G,obs}}
\end{bmatrix}
=
\begin{bmatrix}
d_t \\
1 & 1
\end{bmatrix}
\begin{bmatrix}
1 \\
\mathbf s_t
\end{bmatrix}
+
\begin{bmatrix}
\varepsilon_t^\pi \\
0
\end{bmatrix},
\qquad
t\in\mathcal Q_4.
\end{equation}

In Q1--Q3, the second observation is missing.


% ------------------------------------------------
\subsection{Model 1: HSA Competition Effect on Slope}
% ------------------------------------------------

The slow-moving competition state determines the time-varying
Phillips-curve slope.

\begin{equation}
\kappa_t
=
\kappa_0+\delta\tilde N_t,
\qquad
\pi_t
=
d_t+\delta x_t\tilde N_t+\varepsilon_t^\pi.
\end{equation}

In Q4,

\begin{equation}
\begin{bmatrix}
\pi_t \\
n_t^{\mathrm{G,obs}}
\end{bmatrix}
=
\begin{bmatrix}
d_t+\delta x_t\tilde N_t \\
\tilde N_t+\hat N_t
\end{bmatrix}
+
\begin{bmatrix}
\varepsilon_t^\pi \\
0
\end{bmatrix}.
\end{equation}


% ------------------------------------------------
\subsection{Model 2: HSA Direct Competition Channel}
% ------------------------------------------------

The stationary competition state directly affects quarterly inflation.

\begin{equation}
\pi_t
=
d_t+\theta_0\hat N_t+\varepsilon_t^\pi,
\qquad
\delta=\gamma=0.
\end{equation}

In Q4,

\begin{equation}
\begin{bmatrix}
\pi_t \\
n_t^{\mathrm{G,obs}}
\end{bmatrix}
=
\begin{bmatrix}
d_t+\theta_0\hat N_t \\
\tilde N_t+\hat N_t
\end{bmatrix}
+
\begin{bmatrix}
\varepsilon_t^\pi \\
0
\end{bmatrix}.
\end{equation}

In Q1--Q3, both establishment growth and quarterly inflation provide
information about the unobserved competition state.


% ------------------------------------------------
\subsection{Model 3: HSA Competition Dynamic Effect}
% ------------------------------------------------

The direct competition loading varies with the slow-moving competition
environment.

\begin{equation}
\theta_t
=
\theta_0+\gamma\tilde N_t,
\qquad
\delta=0.
\end{equation}

\begin{equation}
\pi_t
=
d_t
+
\theta_0\hat N_t
+
\gamma\tilde N_t\hat N_t
+
\varepsilon_t^\pi.
\end{equation}

In Q4,

\begin{equation}
\begin{bmatrix}
\pi_t \\
n_t^{\mathrm{G,obs}}
\end{bmatrix}
=
\begin{bmatrix}
d_t+\theta_0\hat N_t+\gamma\tilde N_t\hat N_t \\
\tilde N_t+\hat N_t
\end{bmatrix}
+
\begin{bmatrix}
\varepsilon_t^\pi \\
0
\end{bmatrix}.
\end{equation}

The interaction $\tilde N_t\hat N_t$ makes the inflation measurement equation
bilinear in the latent states.


% ------------------------------------------------
\subsection{Model 4: HSA Joint Competition Effect on Slope and Direct Competition Channel}
% ------------------------------------------------

The full specification combines the competition-dependent slope and dynamic
competition channels.

\begin{equation}
\kappa_t
=
\kappa_0+\delta\tilde N_t,
\qquad
\theta_t
=
\theta_0+\gamma\tilde N_t.
\end{equation}

\begin{equation}
\pi_t
=
d_t
+
\delta x_t\tilde N_t
+
\theta_0\hat N_t
+
\gamma\tilde N_t\hat N_t
+
\varepsilon_t^\pi.
\end{equation}

In Q4, the complete measurement system is

\begin{equation}
\begin{bmatrix}
\pi_t \\
n_t^{\mathrm{G,obs}}
\end{bmatrix}
=
\begin{bmatrix}
d_t
+\delta x_t\tilde N_t
+\theta_0\hat N_t
+\gamma\tilde N_t\hat N_t
\\
\tilde N_t+\hat N_t
\end{bmatrix}
+
\begin{bmatrix}
\varepsilon_t^\pi \\
0
\end{bmatrix}.
\end{equation}

In Q1--Q3,

\begin{equation}
n_t^{\mathrm{G,obs}}=\text{missing},
\qquad
\hat N_t
=
\rho_N\hat N_{t-1}
+
\lambda_E g_t^E
+
\eta_t^{\hat N}.
\end{equation}

The latent quarterly competition path therefore combines three sources of
information: persistence in competition, quarterly business-demography
movements, and the inflation equation, while every observed Q4 value remains
exactly anchored to the original annual Gustavo series.

\begin{equation}
C_t^{\mathrm{slope}}=\delta\tilde N_t x_t,
\qquad
C_t^{\mathrm{direct}}=\theta_0\hat N_t,
\qquad
C_t^{\mathrm{interaction}}=\gamma\tilde N_t\hat N_t.
\end{equation}

% ================================================================
\section{How to Estimate the Model}
% ================================================================

The estimation strategy is Bayesian. The computational method is organized
according to the structure of each model:

\begin{equation}
\text{CES}
\;:\;
\text{Gibbs Sampling},
\qquad
\text{HSA Models 1--2}
\;:\;
\text{Gibbs Sampling + FFBS},
\end{equation}

\begin{equation}
\text{HSA Models 3--4}
\;:\;
\text{Gibbs Sampling + Blocked Conditional FFBS}.
\end{equation}

This strategy applies to all four data cases. Case 4 differs only in its
state transition equation, which additionally uses quarterly establishment
growth as an observed transition input.

% ------------------------------------------------
\subsection{CES Benchmark: Gibbs Sampling}
% ------------------------------------------------

The CES benchmark is

\begin{equation}
\pi_t
=
\alpha E_t\pi_{t+1}
+
\kappa_0x_t
+
\varepsilon_t^\pi,
\qquad
\varepsilon_t^\pi
\sim
\mathcal N(0,\sigma_\pi^2).
\end{equation}

Define

\begin{equation}
\boldsymbol{\beta}_0
=
\begin{bmatrix}
\alpha\\
\kappa_0
\end{bmatrix},
\qquad
\mathbf z_t
=
\begin{bmatrix}
E_t\pi_{t+1}\\
x_t
\end{bmatrix}.
\end{equation}

Stacking over $t$ gives

\begin{equation}
\boldsymbol{\pi}
=
\mathbf Z\boldsymbol{\beta}_0
+
\boldsymbol{\varepsilon}^{\pi},
\qquad
\boldsymbol{\varepsilon}^{\pi}
\sim
\mathcal N
\left(
\mathbf 0,
\sigma_\pi^2\mathbf I
\right).
\end{equation}

With

\begin{equation}
\boldsymbol{\beta}_0
\sim
\mathcal N(\mathbf b_0,\mathbf V_0),
\qquad
\sigma_\pi^2
\sim
\operatorname{IG}(a_{\pi,0},b_{\pi,0}),
\end{equation}

the posterior coefficient update is

\begin{equation}
\mathbf V_n^{-1}
=
\mathbf V_0^{-1}
+
\frac{1}{\sigma_\pi^2}
\mathbf Z^\prime\mathbf Z,
\end{equation}

\begin{equation}
\mathbf b_n
=
\mathbf V_n
\left(
\mathbf V_0^{-1}\mathbf b_0
+
\frac{1}{\sigma_\pi^2}
\mathbf Z^\prime\boldsymbol{\pi}
\right),
\end{equation}

\begin{equation}
\boldsymbol{\beta}_0\mid\cdot
\sim
\mathcal N(\mathbf b_n,\mathbf V_n).
\end{equation}

The variance update is

\begin{equation}
\sigma_\pi^2\mid\cdot
\sim
\operatorname{IG}
\left(
a_{\pi,0}+\frac{T}{2},
\;
b_{\pi,0}
+
\frac12
(\boldsymbol{\pi}-\mathbf Z\boldsymbol{\beta}_0)^\prime
(\boldsymbol{\pi}-\mathbf Z\boldsymbol{\beta}_0)
\right).
\end{equation}


% ------------------------------------------------
\subsection{HSA Models 1--2: Gibbs Sampling with FFBS}
% ------------------------------------------------

For Cases 1--3, the latent competition state is

\begin{equation}
\mathbf s_t
=
\begin{bmatrix}
\tilde N_t\\
\hat N_t
\end{bmatrix},
\end{equation}

with transition equation

\begin{equation}
\mathbf s_t
=
\mathbf F\mathbf s_{t-1}
+
\boldsymbol{\eta}_t,
\qquad
\mathbf F
=
\begin{bmatrix}
1&0\\
0&\rho_N
\end{bmatrix}.
\end{equation}

\begin{equation}
\boldsymbol{\eta}_t
\sim
\mathcal N(\mathbf 0,\mathbf Q),
\qquad
\mathbf Q
=
\begin{bmatrix}
\sigma_{\bar N}^2&0\\
0&\sigma_{\hat N}^2
\end{bmatrix}.
\end{equation}

For Case 4, establishment growth enters the state transition as an observed
input:

\begin{equation}
\mathbf s_t
=
\mathbf F\mathbf s_{t-1}
+
\mathbf B g_t^E
+
\boldsymbol{\eta}_t,
\qquad
\mathbf B
=
\begin{bmatrix}
0\\
\lambda_E
\end{bmatrix}.
\end{equation}

Thus,

\begin{equation}
\tilde N_t
=
\tilde N_{t-1}
+
\eta_t^{\bar N},
\qquad
\hat N_t
=
\rho_N\hat N_{t-1}
+
\lambda_Eg_t^E
+
\eta_t^{\hat N}.
\end{equation}

Models 1 and 2 are conditionally linear Gaussian.

For Model 1,

\begin{equation}
\pi_t
=
\alpha E_t\pi_{t+1}
+
\kappa_0x_t
+
\delta x_t\tilde N_t
+
\varepsilon_t^\pi.
\end{equation}

For Model 2,

\begin{equation}
\pi_t
=
\alpha E_t\pi_{t+1}
+
\kappa_0x_t
+
\theta_0\hat N_t
+
\varepsilon_t^\pi.
\end{equation}

The complete state path is therefore drawn jointly by FFBS.

% ------------------------------------------------
\subsubsection{Forward Filtering}
% ------------------------------------------------

Write the conditional linear-Gaussian system as

\begin{equation}
\mathbf s_t
=
\mathbf c_t
+
\mathbf F_t\mathbf s_{t-1}
+
\mathbf w_t,
\qquad
\mathbf y_t
=
\mathbf d_t
+
\mathbf H_t\mathbf s_t
+
\mathbf v_t,
\end{equation}

with

\begin{equation}
\mathbf w_t\sim\mathcal N(\mathbf 0,\mathbf Q_t),
\qquad
\mathbf v_t\sim\mathcal N(\mathbf 0,\mathbf R_t).
\end{equation}

If

\begin{equation}
\mathbf s_{t-1}\mid\mathbf y_{1:t-1}
\sim
\mathcal N(\mathbf m_{t-1},\mathbf C_{t-1}),
\end{equation}

then the prediction step is

\begin{equation}
\mathbf a_t
=
\mathbf c_t+\mathbf F_t\mathbf m_{t-1},
\qquad
\mathbf P_t
=
\mathbf F_t\mathbf C_{t-1}\mathbf F_t^\prime+\mathbf Q_t.
\end{equation}

The innovation is

\begin{equation}
\mathbf v_t^y
=
\mathbf y_t-\mathbf d_t-\mathbf H_t\mathbf a_t,
\qquad
\mathbf S_t^y
=
\mathbf H_t\mathbf P_t\mathbf H_t^\prime+\mathbf R_t.
\end{equation}

The Kalman gain and filtering update are

\begin{equation}
\mathbf K_t
=
\mathbf P_t\mathbf H_t^\prime
(\mathbf S_t^y)^{-1},
\end{equation}

\begin{equation}
\mathbf m_t
=
\mathbf a_t+\mathbf K_t\mathbf v_t^y,
\qquad
\mathbf C_t
=
\mathbf P_t
-
\mathbf K_t\mathbf S_t^y\mathbf K_t^\prime.
\end{equation}

For mixed-frequency Cases 3 and 4, the competition observation is missing in
Q1--Q3. The missing measurement row is simply omitted from the Kalman update.
No quarterly competition observation is interpolated.

% ------------------------------------------------
\subsubsection{Backward Simulation}
% ------------------------------------------------

At the final period,

\begin{equation}
\mathbf s_T
\sim
\mathcal N(\mathbf m_T,\mathbf C_T).
\end{equation}

For $t=T-1,\ldots,1$,

\begin{equation}
\mathbf J_t
=
\mathbf C_t\mathbf F_{t+1}^\prime
\mathbf P_{t+1}^{-1}.
\end{equation}

Then

\begin{equation}
\mathbf s_t
\mid
\mathbf s_{t+1},\mathbf y_{1:T}
\sim
\mathcal N
\left(
\mathbf m_t
+
\mathbf J_t
(\mathbf s_{t+1}-\mathbf a_{t+1}),
\;
\mathbf C_t
-
\mathbf J_t\mathbf P_{t+1}\mathbf J_t^\prime
\right).
\end{equation}


% ------------------------------------------------
\subsection{HSA Models 3--4: Gibbs Sampling with Blocked Conditional FFBS}
% ------------------------------------------------

Models 3 and 4 contain the bilinear term

\begin{equation}
\gamma\tilde N_t\hat N_t.
\end{equation}

The complete two-dimensional system is therefore not jointly linear in the
latent states. Conditional on either latent state path, however, the other
state enters linearly.

For Model 3,

\begin{equation}
\pi_t
=
\alpha E_t\pi_{t+1}
+
\kappa_0x_t
+
\theta_0\hat N_t
+
\gamma\tilde N_t\hat N_t
+
\varepsilon_t^\pi.
\end{equation}

Conditional on $\hat N_{1:T}$,

\begin{equation}
\pi_t
-
\alpha E_t\pi_{t+1}
-
\kappa_0x_t
-
\theta_0\hat N_t
=
\gamma\hat N_t\tilde N_t
+
\varepsilon_t^\pi.
\end{equation}

Conditional on $\tilde N_{1:T}$,

\begin{equation}
\pi_t
-
\alpha E_t\pi_{t+1}
-
\kappa_0x_t
=
(\theta_0+\gamma\tilde N_t)\hat N_t
+
\varepsilon_t^\pi.
\end{equation}

For Model 4,

\begin{equation}
\pi_t
=
\alpha E_t\pi_{t+1}
+
\kappa_0x_t
+
\delta x_t\tilde N_t
+
\theta_0\hat N_t
+
\gamma\tilde N_t\hat N_t
+
\varepsilon_t^\pi.
\end{equation}

Conditional on $\hat N_{1:T}$,

\begin{equation}
\pi_t
-
\alpha E_t\pi_{t+1}
-
\kappa_0x_t
-
\theta_0\hat N_t
=
(\delta x_t+\gamma\hat N_t)\tilde N_t
+
\varepsilon_t^\pi.
\end{equation}

Conditional on $\tilde N_{1:T}$,

\begin{equation}
\pi_t
-
\alpha E_t\pi_{t+1}
-
\kappa_0x_t
-
\delta x_t\tilde N_t
=
(\theta_0+\gamma\tilde N_t)\hat N_t
+
\varepsilon_t^\pi.
\end{equation}

Hence the Gibbs sampler alternates between two exact conditional FFBS draws:

\begin{equation}
\tilde N_{1:T}^{(g)}
\sim
p
\left(
\tilde N_{1:T}
\mid
\hat N_{1:T}^{(g-1)},
\boldsymbol{\Theta}^{(g-1)},
\mathcal D
\right),
\end{equation}

\begin{equation}
\hat N_{1:T}^{(g)}
\sim
p
\left(
\hat N_{1:T}
\mid
\tilde N_{1:T}^{(g)},
\boldsymbol{\Theta}^{(g-1)},
\mathcal D
\right).
\end{equation}

The same blocked conditional FFBS strategy is used in Case 4. The
establishment-growth term simply enters the conditional transition mean of
$\hat N_t$.


% ------------------------------------------------
\subsection{Transition-Parameter Gibbs Updates}
% ------------------------------------------------

For Cases 1--3,

\begin{equation}
\hat N_t
=
\rho_N\hat N_{t-1}
+
\eta_t^{\hat N}.
\end{equation}

For Case 4,

\begin{equation}
\hat N_t
=
\rho_N\hat N_{t-1}
+
\lambda_E g_t^E
+
\eta_t^{\hat N}.
\end{equation}

Thus Case 4 can be written as the Gaussian regression

\begin{equation}
\mathbf r_E
=
\mathbf X_E
\begin{bmatrix}
\rho_N\\
\lambda_E
\end{bmatrix}
+
\boldsymbol{\eta}^{\hat N},
\end{equation}

where

\begin{equation}
\mathbf r_E
=
\begin{bmatrix}
\hat N_2\\
\vdots\\
\hat N_T
\end{bmatrix},
\qquad
\mathbf X_E
=
\begin{bmatrix}
\hat N_1 & g_2^E\\
\vdots & \vdots\\
\hat N_{T-1} & g_T^E
\end{bmatrix}.
\end{equation}

With a Gaussian prior, $(\rho_N,\lambda_E)$ has a Gaussian conditional
posterior, truncated only to enforce $|\rho_N|<1$.

The state residuals are

\begin{equation}
u_t^{\bar N}
=
\tilde N_t-\tilde N_{t-1},
\end{equation}

and

\begin{equation}
u_t^{\hat N}
=
\begin{cases}
\hat N_t-\rho_N\hat N_{t-1},
& \text{Cases 1--3},\\[4pt]
\hat N_t-\rho_N\hat N_{t-1}-\lambda_Eg_t^E,
& \text{Case 4}.
\end{cases}
\end{equation}

The corresponding innovation variances are updated from inverse-gamma full
conditionals.


% ------------------------------------------------
\subsection{Complete Gibbs Sweep}
% ------------------------------------------------

For HSA Models 1 and 2:

\begin{enumerate}
\item Draw the complete state path $\mathbf s_{1:T}$ by FFBS.
\item Draw the inflation coefficients conditional on the state path.
\item Draw $\sigma_\pi^2$.
\item Draw the transition parameters $\rho_N$ and, in Case 4, $\lambda_E$.
\item Draw the state-innovation variances and applicable measurement variance.
\item Store the post-warm-up draw.
\end{enumerate}

For HSA Models 3 and 4:

\begin{enumerate}
\item Draw $\tilde N_{1:T}$ by FFBS conditional on $\hat N_{1:T}$.
\item Draw $\hat N_{1:T}$ by FFBS conditional on $\tilde N_{1:T}$.
\item Draw the inflation coefficients conditional on both state paths.
\item Draw $\sigma_\pi^2$.
\item Draw the transition parameters $\rho_N$ and, in Case 4, $\lambda_E$.
\item Draw the state-innovation variances and applicable measurement variance.
\item Store the post-warm-up draw.
\end{enumerate}


% ------------------------------------------------
\subsection{Convergence and Computational Validation}
% ------------------------------------------------

All models are estimated using multiple dispersed chains. Convergence is
assessed using

\begin{equation}
\widehat R,
\qquad
ESS_{\mathrm{bulk}},
\qquad
ESS_{\mathrm{tail}},
\end{equation}

together with trace plots and autocorrelation summaries.

For Models 1 and 2, the full two-dimensional FFBS is used directly. For
Models 3 and 4, the blocked conditional FFBS implementation is validated
against the linear special case obtained by setting $\gamma=0$.

For mixed-frequency Cases 3 and 4, Q1--Q3 remain missing throughout the
estimation. The quarterly latent competition path is inferred by the
state-space smoother rather than constructed by deterministic interpolation.

\end{document}
